\section{星眠与 Dokie：两条产品线}

前面解释了 all in 的创业判断，现在需要看这个判断具体落在哪些产品上。星眠和 Dokie 分别代表陪伴内容和生产力 Agent 两条路线。

沐言智语现在有两条主线：AI 乙女游戏星眠，以及 Agent 产品 Dokie。前者延续张月光对 Live2D、陪伴和女性向内容的长期兴趣；后者是他探索一年后决定 all in 的生产力方向。

\begin{figure}[H]
\centering
\includegraphics[width=0.94\textwidth]{xingmian-dokie-two-bets.png}
\caption{沐言智语的两条产品线：星眠与 Dokie。}
\end{figure}

\begin{knowledgebox}{读图：一个做内容，一个做能力边界}
星眠验证 AI 内容和陪伴，重点是内容、角色、频次和交互；Dokie 从 PPT/创作表达切入，目标是突破用户能力边界。二者都不是单纯炫技，而是在寻找用户会长期交互、长期记住、长期付费的入口。
\end{knowledgebox}

\subsection{AI 游戏为什么窄}

张月光认为，AI 游戏的可行空间比很多人想得窄。实时生成视频、无限互动影游、开放世界动态剧情，今天技术还很难稳定落地。AI 最容易加分的，是对话陪伴和角色关系，因此他更看好 AI 乙女游戏，而不是泛泛的 AI 游戏。

\begin{warningbox}{AI 游戏不是所有游戏加 AI}
游戏本质仍是内容产业。AI 可以改变某些内容体验，但不能自动解决剧情、成本、稳定性和用户留存。可行场景往往比宏大愿景更窄。
\end{warningbox}

\subsection{Dokie 为什么从 PPT 切入}

Dokie 从 PPT 切入，但张月光反复强调，它不会永远是 PPT 工具。PPT 是一个创作与表达的入口，真正目标是“帮我突破能力边界的 AI 朋友”。这解释了为什么他愿意 all in Agent，而不愿把全部资金押在此前那些小探索上。

\subsection{本章小结}

星眠和 Dokie 分别代表 AI 人口的两种形态：陪伴/内容中的 AI 个体，以及生产力/创作中的 AI 个体。前者要控制内容成本和陪伴边界，后者要证明能力边界真的被突破。

